% Web source: docs/05statistical_mechanics/01/01_030.md
\addcontentsline{toc}{section}{1-3　スターリングの公式}

\begin{mondai}{1-3}{スターリングの公式}{基本}
統計力学において，構成粒子数がアボガドロ数規模に達する多体系を扱う際，階乗 $N!$ の計算が頻出する。この巨大な階乗の対数を扱いやすい形に変形するため，スターリングの公式 $\log N! \simeq N \log N - N$ の導出とその適用について，以下の問いに答えよ。ただし，対数は自然対数とする。

\begin{enumerate}
  \item 関数 $\log x$ が単調増加することを用いて，階乗の対数 $\log N! = \sum_{x=1}^N \log x$ が次の不等式を満たすことを示せ。
  \begin{equation*}
  \int_1^N \log x \, dx \le \log N! \le \int_1^N \log x \, dx + \log N
  \end{equation*}
  さらに，$N$ が十分に大きいとき，右辺に現れる補正項 $\log N$ が主要項に比べて無視できることを示せ。

  \item 部分積分法を用いることで定積分 $\int_1^N \log x \, dx$ を計算し，スターリングの公式 $\log N! \simeq N \log N - N$ を導出せよ。

  \item 導出した公式を用いて，独立な $N$ 個の二状態系においてちょうど半数の $N/2$ 個の粒子が励起状態にあるときの微視的状態数
  \begin{equation*}
  \Omega = \frac{N!}{\left(\frac{N}{2}\right)! \left(\frac{N}{2}\right)!}
  \end{equation*}
  の対数 $\log \Omega$ を，$N$ を用いた簡潔な形に近似せよ。
\end{enumerate}
\end{mondai}

\solutionheading
\begin{solutionparts}

\solutionpart{(1)}
階乗の対数は，対数の性質 $\log(ab) = \log a + \log b$ を用いることで，以下のように和の形で書き下すことができます。
\begin{equation*}
\log N! = \log (1 \times 2 \times \cdots \times N) = \sum_{x=1}^N \log x
\end{equation*}
$\log x$ は単調増加関数です。区間 $[k-1,k]$ では $\log x \le \log k$ であるため，
\begin{equation*}
\int_{k-1}^{k}\log x\,dx \le \log k
\end{equation*}
が成り立ちます。$k=2$ から $N$ まで足し合わせると，$\log 1=0$ より
\begin{equation*}
\int_1^N\log x\,dx \le \sum_{k=1}^N\log k=\log N!
\end{equation*}
を得ます。一方，区間 $[k,k+1]$ では $\log k\le\log x$ なので，$k=1$ から $N-1$ まで足し合わせると
\begin{equation*}
\log N! - \log N \le \int_1^N\log x\,dx
\end{equation*}
となります。したがって，
\begin{equation*}
\int_1^N \log x \, dx \le \log N! \le \int_1^N \log x \, dx + \log N
\end{equation*}
が成り立ちます。積分の主要項は $N\log N$ のオーダーであるのに対し，補正項は $\log N$ のオーダーです。その比は $1/N$ で0に近づくため，$N$ が十分に大きければ和を積分で近似できます。

\solutionpart{(2)}
$y = \log x$ の積分を部分積分法を用いて計算します。$\log x$ の前に隠れている $1 = (x)^\prime$ を意識して計算を進めます。
\begin{equation*}
\int_1^N \log x \, dx = \left[ x \log x \right]_1^N - \int_1^N x \cdot \frac{1}{x} \, dx
\end{equation*}
それぞれの項を評価します。第1項に上限 $N$ と下限 $1$ を代入すると，$\log 1 = 0$ であることから $N \log N$ となります。第2項の被積分関数は $1$ になるため，積分結果は $N - 1$ です。
\begin{equation*}
\int_1^N \log x \, dx = N \log N - (N - 1) = N \log N - N + 1
\end{equation*}
ここで $N$ が非常に大きい極限を考えているため，定数項の $1$ は $N$ や $N \log N$ に比べて無視することができます。これにより，スターリングの公式が導かれます。
\begin{equation*}
\log N! \simeq N \log N - N
\end{equation*}

\solutionpart{(3)}
状態数 $\Omega$ の対数をとり，対数の性質を用いて分子と分母を分解します。
\begin{equation*}
\log \Omega = \log N! - 2 \log \left(\frac{N}{2}\right)!
\end{equation*}
ここに，先ほど導出したスターリングの公式をそれぞれ適用します。
\begin{equation*}
\log \Omega \simeq (N \log N - N) - 2 \left[ \frac{N}{2} \log \left(\frac{N}{2}\right) - \frac{N}{2} \right]
\end{equation*}
右辺の第2項の括弧を展開して整理します。
\begin{equation*}
\log \Omega \simeq N \log N - N - N \left( \log N - \log 2 \right) + N
\end{equation*}
\begin{equation*}
\log \Omega \simeq N \log N - N - N \log N + N \log 2 + N
\end{equation*}
共通する項を相殺することで，以下の簡潔な結果が得られます。
\begin{equation*}
\log \Omega \simeq N \log 2
\end{equation*}

\end{solutionparts}
